\chapter{Lorentz Force Law}

\section*{Introduction}

The Lorentz force law is the bridge between electricity, magnetism, and mechanics. It tells us exactly how an electric or magnetic field pushes on a charged particle, and in doing so it unifies two previously separate forces into a single, elegant expression. First formulated by Hendrik Lorentz in 1892, this law is the operational heart of every particle accelerator, mass spectrometer, and plasma confinement device ever built. It completes the classical description of electromagnetism: Maxwell's equations tell you what fields exist; the Lorentz force law tells you what those fields \emph{do} to matter.

Imagine a single charged particle---say, an electron---moving through space where electric and magnetic fields coexist. The electric field exerts a force along its own direction, either accelerating or decelerating the charge. The magnetic field, by contrast, never speeds a particle up or slows it down; it only bends its path. The magnetic force is always perpendicular to the velocity, which means it does zero work and changes the particle's direction without changing its speed. The total force is the vector sum of these two effects, and it is this total force that governs everything from the deflection of cosmic rays in Earth's magnetosphere to the confinement of hundred-million-degree plasma in a tokamak fusion reactor.
\section*{Fundamental Equation Used}

\begin{equation}
\mathbf{F} = q\!\left(\mathbf{E} + \mathbf{v} \times \mathbf{B}\right)
\end{equation}
\section*{Assumptions}

\textbf{Assumptions:} The particle is a point charge with well-defined position and velocity. The fields $\mathbf{E}$ and $\mathbf{B}$ are the total (external plus self) fields at the particle's location. The equation is relativistically exact as written; no approximation $v \ll c$ is required in the force expression itself.


\textbf{Limitations:} Does not account for radiation reaction (the force a charged particle exerts on itself via its own emitted radiation). Breaks down at quantum scales where the classical trajectory is invalid (replaced by quantum electrodynamics). Spin--magnetic-field interactions require the additional Pauli term $-\boldsymbol{\mu} \cdot \mathbf{B}$.


\section*{Mathematical Derivation}

\textbf{Step 1: Start from the electric force on a charge.}

A point charge $q$ at position $\mathbf{r}$ in an electric field $\mathbf{E}(\mathbf{r},t)$ experiences a force
\begin{equation}
\mathbf{F}_E = q\,\mathbf{E}
\end{equation}
This is the empirical law of Coulomb and the defining relation of the electric field. The force is proportional to $q$ and parallel to $\mathbf{E}$.

\textbf{Step 2: Consider a charge moving in a magnetic field in its rest frame.}

In the rest frame of the charge, a purely magnetic field $\mathbf{B}$ produces no force because $\mathbf{v} = 0$. However, if we boost to a frame where the charge moves with velocity $\mathbf{v}$, the magnetic field transforms into a mixture of electric and magnetic fields. In the rest frame, the Lorentz transformation gives a motional electric field $\mathbf{E}' = \gamma(\mathbf{v}\times\mathbf{B})$. At non-relativistic speeds ($\gamma \approx 1$), this becomes $\mathbf{E}' \approx \mathbf{v}\times\mathbf{B}$.

\textbf{Step 3: Write the force in the lab frame via Lorentz transformation.}

The force on a charge transforms under Lorentz boosts. Starting from the covariant expression, the spatial part of the force transforms as:
\begin{equation}
\mathbf{F} = q(\mathbf{E} + \mathbf{v}\times\mathbf{B})
\end{equation}
This can be verified by applying the Lorentz force law $f^\mu = qF^{\mu\nu}u_\nu$ and extracting the spatial components. The key point is that the cross-product $\mathbf{v}\times\mathbf{B}$ arises naturally from the antisymmetric structure of the electromagnetic field tensor.

\textbf{Step 4: Verify consistency with Maxwell's equations.}

The Lorentz force law is consistent with Maxwell's equations in the following sense: if we have a distribution of charges and currents described by $\rho$ and $\mathbf{J}$, then Maxwell's equations determine the fields, and the Lorentz force law determines how those fields push the charges. Together, they form a closed system:
\begin{align}
\nabla\cdot\mathbf{E} &= \frac{\rho}{\epsilon_0} \\
\nabla\times\mathbf{B} &= \mu_0\mathbf{J} + \mu_0\epsilon_0\frac{\partial\mathbf{E}}{\partial t} \\
\mathbf{F} &= q(\mathbf{E} + \mathbf{v}\times\mathbf{B})
\end{align}

\textbf{Step 5: Write the non-relativistic equations of motion.}

Applying Newton's second law ($\mathbf{F} = m\,d\mathbf{v}/dt$), the Lorentz force gives the equation of motion for a charged particle:
\begin{equation}
m\frac{d\mathbf{v}}{dt} = q(\mathbf{E} + \mathbf{v}\times\mathbf{B})
\end{equation}
For the relativistic case, replace $m\,d\mathbf{v}/dt$ with $d\mathbf{p}/dt$ where $\mathbf{p} = \gamma m\mathbf{v}$. The equation remains valid at all speeds when written in this form.

\textbf{Step 6: Verify the magnetic force does no work.}

The power delivered by the Lorentz force is:
\begin{equation}
P = \mathbf{F}\cdot\mathbf{v} = q\mathbf{E}\cdot\mathbf{v} + q(\mathbf{v}\times\mathbf{B})\cdot\mathbf{v}
\end{equation}
Since $(\mathbf{v}\times\mathbf{B})\cdot\mathbf{v} = 0$ (the cross product is perpendicular to $\mathbf{v}$), only the electric field does work:
\begin{equation}
P = q\,\mathbf{E}\cdot\mathbf{v}
\end{equation}
This confirms that magnetic fields can only redirect charged particles, not accelerate them.

\textbf{Step 7: Write the fully covariant form.}

The Lorentz force law can be expressed covariantly using the electromagnetic field tensor $F^{\mu\nu}$ and the four-velocity $u^\nu = \gamma(c,\mathbf{v})$:
\begin{equation}
f^\mu = q\,F^{\mu\nu}u_\nu
\end{equation}
where
\begin{equation}
F^{\mu\nu} = \begin{pmatrix} 0 & -E_x/c & -E_y/c & -E_z/c \\ E_x/c & 0 & -B_z & B_y \\ E_y/c & B_z & 0 & -B_x \\ E_z/c & -B_y & B_x & 0 \end{pmatrix}
\end{equation}
This is the most fundamental form of the Lorentz force law, valid in all frames of reference.

\section*{Characteristics of the Equation}

The Lorentz force is a vector equation in three dimensions. The electric contribution $q\mathbf{E}$ is parallel (or antiparallel, for negative charges) to $\mathbf{E}$. The magnetic contribution $q\mathbf{v}\times\mathbf{B}$ is perpendicular to both $\mathbf{v}$ and $\mathbf{B}$, forming a right-handed triplet. The magnitude of the magnetic force is $|q|vB\sin\phi$, where $\phi$ is the angle between $\mathbf{v}$ and $\mathbf{B}$. The magnetic force vanishes entirely when the velocity is parallel to the field. The electric force can do work; the magnetic force cannot. This asymmetry is profound: it means magnetic fields can steer particles but never heat them directly, while electric fields can do both.
\section*{Solution Methods}

In the simplest case---a uniform magnetic field with no electric field---the particle executes helical motion. Decompose the velocity into components parallel and perpendicular to $\mathbf{B}$. The parallel component is unaffected; the perpendicular component produces uniform circular motion with cyclotron frequency $\omega_c = qB/m$ and radius $r = mv_\perp / (qB)$, the Larmor radius. When both fields are present and configured so that $qE = qvB$ (i.e., $v = E/B$), the electric and magnetic forces exactly cancel for particles with that specific speed---this is the principle of the velocity selector. More complex configurations (non-uniform fields, time-varying fields) require numerical integration of the equations of motion, often via the Boris pusher algorithm in plasma physics.
\section*{Verifications}

The Lorentz force law can be verified experimentally by measuring the deflection of an electron beam in known electric and magnetic fields and confirming that the observed trajectory matches $\mathbf{F} = q(\mathbf{E} + \mathbf{v}\times\mathbf{B})$. J.J. Thomson's 1897 measurement of the electron's charge-to-mass ratio using crossed fields is a direct confirmation. The Hall effect provides a macroscopic verification: the measured Hall voltage $V_H = IB/(nqt)$ (for a thin slab) agrees precisely with the Lorentz force prediction. In particle physics, the helical tracks of charged particles in bubble chambers and cloud chambers are visual confirmations of the magnetic force component.
\section*{Applications}

The Lorentz force is the operating principle of cyclotrons and synchrotrons (particle accelerators that use magnetic fields to bend particle paths and electric fields to accelerate them), mass spectrometers (which separate ions by their charge-to-mass ratio via differential bending in magnetic fields), magnetohydrodynamic (MHD) generators and electromagnetic flow meters (which exploit the Hall voltage created when a conducting fluid moves through a magnetic field), and the Hall effect itself---the development of a transverse voltage in a current-carrying conductor placed in a magnetic field. It also governs the motion of charged particles in Earth's radiation belts, solar wind interactions with planetary magnetospheres, and magnetic confinement schemes for controlled fusion.
\section*{What Richard Feynman Would Have Asked}

\textit{``If I have a charged particle moving through a region with both electric and magnetic fields, what is really happening to that particle at the most fundamental level---am I seeing two separate forces, or one?''}

This is the kind of question Feynman loved, because it forces us to confront the artificiality of our labels. There is no such thing as a ``purely electric'' or ``purely magnetic'' force---these are descriptions that depend on your frame of reference. What one observer calls a purely electric field, another observer moving relative to the first will describe as a mixture of electric and magnetic fields. The quantity that is truly frame-independent is the electromagnetic field tensor $F^{\mu\nu}$, an antisymmetric $4\times 4$ matrix that encodes all six components of $\mathbf{E}$ and $\mathbf{B}$. The Lorentz force is the response of a charge to this single, unified object. The decomposition into $\mathbf{E}$ and $\mathbf{B}$ is our convenience, not nature's.

\textit{``What would happen to the magnetic force if I could somehow make the charge move faster and faster, approaching the speed of light?''}

In classical electromagnetism, the Lorentz force law as written does not break down as $v\to c$. The equation $\mathbf{F} = q(\mathbf{E} + \mathbf{v}\times\mathbf{B})$ remains valid at all speeds, but you must be careful about what $\mathbf{F}$ means: it is the rate of change of the relativistic momentum $\mathbf{p} = \gamma m\mathbf{v}$, not $m\mathbf{a}$. So as $v\to c$, the same force produces less and less acceleration because the inertia grows without bound. In a particle accelerator, this is exactly what happens: the magnets must become stronger and stronger to bend the path of a particle that is becoming increasingly resistant to changes in its velocity. The 7 TeV protons in the LHC are deflected by fields that would be absurdly strong for a non-relativistic particle, yet they barely curve because $\gamma \approx 7500$.

\textit{``If magnetic fields do no work, how can a generator produce electricity? Doesn't the magnetic field have to do work on the charges?''}

This is a beautifully subtle point. In a generator, the magnetic field does indeed exert no work on the moving charges---the magnetic force is always perpendicular to their velocity. But the generator is being \emph{mechanically driven}: an external agent (a turbine, a hand crank) pushes the conductor through the magnetic field, and it is this mechanical agent that does the work. The magnetic field acts as an intermediary, redirecting the mechanical force into an electromotive force (emf) along the conductor. The electrons are pushed sideways by the magnetic force, accumulating on one end of the wire and creating an electric field; it is this induced electric field that ultimately drives the current and delivers electrical energy to the load. The energy budget closes perfectly: mechanical work in equals electrical work out (minus losses), and the magnetic field stores and transfers energy without consuming any.

\textit{``Is the cyclotron frequency truly independent of the particle's speed, or is that an approximation?''}

For a non-relativistic charged particle in a uniform magnetic field, the cyclotron frequency $\omega_c = qB/m$ is exactly independent of speed---a larger orbit just means the particle travels a longer circumference in the same time. This is the principle that made cyclotrons possible: you can accelerate particles in spiral orbits with a fixed-frequency alternating voltage. However, as the particle becomes relativistic, its effective mass increases as $\gamma m$, and the cyclotron frequency drops: $\omega_c = qB/(\gamma m)$. The synchrotron was invented to address this---it modulates the driving frequency (or the magnetic field) to track the changing $\gamma$. In extreme cases (synchrotron light sources, the LHC), the frequency modulation is substantial, and the simple constant-frequency cyclotron is limited to energies where $\gamma \approx 1$.

\textit{``Could you use the Lorentz force to build a perfect velocity selector, or are there fundamental limits?''}

A velocity selector using crossed $\mathbf{E}$ and $\mathbf{B}$ fields is theoretically perfect in the classical limit: particles with $v = E/B$ pass through undeflected, while all others are deflected. The selectivity depends only on how well you can control and uniformize $\mathbf{E}$ and $\mathbf{B}$ and how narrow you can make the exit slit. In practice, the limits come from fringe fields, field inhomogeneities, and the finite width of the beam. At the quantum level, there is a more fundamental issue: even a perfectly selected particle has a position-momentum uncertainty that gives it a spread in both transverse position and transverse momentum, so the ``selected'' velocity is not infinitely precise. For most engineering purposes, though, this quantum limit is negligible, and velocity selectors achieve precisions of parts in $10^4$ or better with careful field design.

\bigskip
\hrule
\bigskip
%=============================================================================
\section*{FAQ}

\textit{Q: Why does the magnetic force do no work?}
A: Because it is always perpendicular to the velocity. Work is defined as $\mathbf{F}\cdot d\mathbf{r} = \mathbf{F}\cdot\mathbf{v}\,dt$, and if $\mathbf{F}\perp\mathbf{v}$, this dot product is identically zero at every instant. The magnetic field can redirect a particle but cannot change its kinetic energy.

\textit{Q: Does the Lorentz force law apply to neutral particles?}
A: Not directly. A neutral particle ($q=0$) experiences no Lorentz force. However, neutral particles can have magnetic dipole moments (like the neutron) and interact with \emph{gradients} of magnetic fields through the force $\mathbf{F} = \nabla(\boldsymbol{\mu}\cdot\mathbf{B})$. This is a separate effect from the Lorentz force.

\textit{Q: Is the Lorentz force law consistent with special relativity?}
A: Yes. The equation $\mathbf{F} = q(\mathbf{E} + \mathbf{v}\times\mathbf{B})$ is the spatial part of the fully covariant expression $f^\mu = qF^{\mu\nu}u_\nu$, where $F^{\mu\nu}$ is the electromagnetic field tensor and $u_\nu$ is the four-velocity. What appears as separate electric and magnetic fields in one frame transforms into a mixture of both in another frame; the Lorentz force law is frame-covariant.
\section*{Important Bibliography}
\begin{enumerate}
\item Griffiths, D.J., \textit{Introduction to Electrodynamics}, 4th ed. (Cambridge University Press, 2017), Chapter 5.
\item Jackson, J.D., \textit{Classical Electrodynamics}, 3rd ed. (Wiley, 1999), Chapter 12.
\item Lorentz, H.A., \textit{The Theory of Electrons} (Teubner, 1909), Chapters 1--2.
\item Purcell, E.M. and Morin, D.J., \textit{Electricity and Magnetism}, 3rd ed. (Cambridge University Press, 2013), Chapter 12.
\item Zangwill, A., \textit{Modern Electrodynamics} (Cambridge University Press, 2013), Chapter 16.
\end{enumerate}


